\documentclass[10pt,a4paper]{amsart}
\usepackage[margin=19mm]{geometry}
\usepackage{amsmath,amssymb,mathtools}
\usepackage[scr=rsfs,cal=euler]{mathalfa}
\usepackage{xfrac}
\usepackage[unicode,hidelinks,bookmarks=false]{hyperref}
\newcommand{\ie}{i.e.\ }
\newcommand{\cf}{cf.\ }
\newcommand{\resp}{respectively\ }
\newcommand{\Subsection}{\subsection}
\providecommand{\nobreakdash}{\nobreak\hbox{-}}
\setlength{\emergencystretch}{2em}
\multlinegap0pt
\usepackage{amssymb}
\usepackage{enumerate}
\usepackage{mathtools}
\usepackage[shortlabels]{enumitem}
\usepackage{xcolor}
\usepackage{cancel}
\usepackage[normalem]{ulem}
\usepackage{tikz}
\newtheorem{thm}{Theorem}
\title{A quantitative container characterization \\of one-sided testability}
\author{Gaia Carenini, Cameron Seth, Yuichi Yoshida}
\date{Source: arXiv:2608.01523v1 (CC BY 4.0)}
\begin{document}
\maketitle

\noindent\emph{Source and licence.} A quantitative container characterization \\of one-sided testability. Version arXiv:2608.01523v1 (CC BY 4.0), distributed by arXiv under the Creative Commons Attribution 4.0 International licence, http://creativecommons.org/licenses/by/4.0/, as recorded in the arXiv licence metadata for this exact version and shown on the abstract page https://arxiv.org/abs/2608.01523v1. Full licence notice: \texttt{licenses/arxiv-2608.01523-CC-BY-4.0.txt}.\par\smallskip

\noindent\textit{Editorial scope.} Counted unit: the article's thm thm:removal-container-equivalence (source0.tex lines 637--656). The article's complete original proof of it is delivered with it (source0.tex lines 658--772, one \texttt{proof} environment of 115 lines). No mathematics is written, repaired or re-derived: the statement and the proof are the article's own lines, in the article's own notation, and no retained line is substituted or re-flowed.  No further source-local context is needed: every cross-reference of every retained line resolves inside the two delivered spans.  Nothing was omitted from any retained span, so the package carries no omission record.  No retained line contains a citation, so the wrapper carries no bibliography and no bibliography entry.  No retained line contains a figure environment, an image-inclusion command or drawing code, so no image is delivered and the package declares no asset dependency.  Editorial formatting only, in addition: an AMS \texttt{amsart} preamble that restates the article's own declarations and macros used by the retained lines, namely the environments \texttt{thm} on separate counters, together with the standard packages those lines need.  The retained environments print in this document as Theorem 1 in order of appearance, whereas the article prints the same items with the numbering of its own full document.  The complete native source of the article as distributed in the arXiv e-print for this exact version is bundled unchanged as \texttt{sources/206638/source0.tex}.
\begin{thm}[Removal--container equivalence]\label{thm:removal-container-equivalence}
Let $\Pi$ be a hereditary graph property.
\begin{enumerate}
\item If $\Pi$ admits an $(\varepsilon,\delta,N,s)$-removal lemma, then it
admits the following ordered container statement: for every integer
$n\geq\max\{N,s\}$ and every $n$-vertex graph $G=(V,E)$ that is
$\varepsilon$-far from $\Pi$, there are maps
\(f:\{U\subseteq V:G[U]\in\Pi\} \longrightarrow \bigcup_{j=0}^{s-1}V_{\neq}^{j}, \qquad C:\operatorname{im}(f) \longrightarrow \binom{V}{\leq(1-\delta/s)n},\)
such that every entry of $f(U)$ belongs to $U$ and
$U\subseteq C(f(U))$ whenever $G[U]\in\Pi$.

\item Conversely, suppose that $\Pi$ admits an ordered container statement
with proximity parameter $\varepsilon$, shrinkage $\eta$, threshold $N$,
and fingerprint bound $q$.  Then, for every integer $t\geq q$ such that
\(\sum_{j=0}^{q}t^j(1-\eta)^{t-j}<1,\)
the property $\Pi$ admits an $(\varepsilon,\delta_t,N,t)$-removal lemma,
where
\(\delta_t=1-\sum_{j=0}^{q}t^j(1-\eta)^{t-j}>0.\)
\end{enumerate}
\end{thm}

\begin{proof}
We first prove the forward implication.  Let $G=(V,E)$ be an $n$-vertex graph
with $n\geq\max\{N,s\}$ that is $\varepsilon$-far from $\Pi$, and let $H$ be the
$s$-uniform hypergraph on $V$ whose edges are the sets $S\in\binom Vs$ for
which $G[S]\notin\Pi$.  By hypothesis,
$e(H)\geq\delta\binom ns$.  Every set $U\subseteq V$ with $G[U]\in\Pi$ is an
independent set of $H$.

Fix a total order on $V$, used only to make all choices canonical.  If $s=1$,
let every independent set have the empty fingerprint and put
$C(())=V\setminus E(H)$, where the edge set of the $1$-uniform hypergraph is
identified with a subset of $V$.  Then
$|C(())|\leq(1-\delta)n$, and the conclusion follows.  We may therefore assume
$s\geq2$.

If $n=s$, then $H$ is nonempty and its unique possible edge is $V$.  Hence
every independent set is a proper subset of $V$.  Assign to an independent set
its increasing ordering as fingerprint and use the set itself as its
container.  Since $\delta\leq1$, every such container has size at most
$n-1\leq(1-\delta/s)n$.

Assume henceforth that $n>s$.  If an independent set $U$ has
$|U|\leq s-2$, let $f(U)$ be its increasing ordering and set $C(f(U))=U$.
These fingerprints have length at most $s-2$.  Now let $|U|\geq s-1$.
Starting from $H^{(s)}=H$, define hypergraphs
$H^{(s)},H^{(s-1)},\ldots,H^{(1)}$ recursively.  For
$i=s,s-1,\ldots,2$, choose a vertex
$v_i\in U\cap V(H^{(i)})$ of maximum degree among the vertices of
$U\cap V(H^{(i)})$, breaking ties by the fixed order, and put
\(X_i=\{u\in V(H^{(i)}): \deg_{H^{(i)}}(u)>\deg_{H^{(i)}}(v_i)\}.\)
Let $H^{(i-1)}$ be the $(i-1)$-uniform link of $v_i$ in $H^{(i)}$, on the
vertex set $V(H^{(i)})\setminus\{v_i\}$.  Since $U$ is independent in
$H^{(i)}$, the set $U\cap V(H^{(i-1)})$ is independent in $H^{(i-1)}$; in
particular, the construction is well defined through the last step.

Set
\(f(U)=(v_s,v_{s-1},\ldots,v_2)
\)
and define
\(C(f(U)) =V\setminus\left(E(H^{(1)})\cup\bigcup_{i=2}^{s}X_i\right).\)
Here again $E(H^{(1)})$ is viewed as a set of vertices.  This is a genuine
function of the fingerprint: starting from $H$, the ordered tuple
$(v_s,\ldots,v_2)$ uniquely reconstructs all links $H^{(i)}$, all sets $X_i$,
and hence the displayed container.  Thus two independent sets with the same
fingerprint receive the same container.  Moreover, fingerprints arising in
this case have length exactly $s-1$, so they cannot coincide with the
fingerprints used for sets of size at most $s-2$.

By maximality of $v_i$ inside $U\cap V(H^{(i)})$, every $X_i$ is disjoint from
$U$.  Also, $U\cap V(H^{(1)})$ avoids $E(H^{(1)})$.  Consequently
$U\subseteq C(f(U))$.

It remains to prove the shrinkage.  Write
$n_i=|V(H^{(i)})|=n-s+i$ and $m_i=e(H^{(i)})$.  Then
$m_{i-1}=\deg_{H^{(i)}}(v_i)$ for $2\leq i\leq s$.  Let $k$ be the least
index in $[s]$ such that either $k=1$ or
\(m_{k-1}<\frac{k-1}{n_k-1}m_k.\)
For every $j=k+1,\ldots,s$ we therefore have
$m_{j-1}\geq\frac{j-1}{n_j-1}m_j$, and hence
\[
m_k
\geq
\left(\prod_{j=k+1}^{s}\frac{j-1}{n_j-1}\right)m_s
\geq
\frac{\delta n}{s}\binom{n_k-1}{k-1}.
\]
If $k=1$, then $|E(H^{(1)})|=m_1\geq\delta n/s$, which gives the required
bound on $|C(f(U))|$.

Suppose that $k\geq2$ and put
\[
Y=\left\{u\in V(H^{(k)}):
\deg_{H^{(k)}}(u)>\frac{k-1}{n_k-1}m_k\right\}.
\]
Since $n_k\geq k+1$, the average degree $km_k/n_k$ is strictly larger than
the threshold in this definition.  Furthermore,
\(|Y|\geq\frac{m_k}{\binom{n_k-1}{k-1}}.\)
Indeed, otherwise the contribution to the degree sum from $Y$ is less than
$m_k$, while the remaining at most $n_k-1$ vertices contribute at most
$(k-1)m_k$, contradicting that the total degree sum is $km_k$.  By the choice
of $k$,
$\deg_{H^{(k)}}(v_k)=m_{k-1}<\frac{k-1}{n_k-1}m_k$, so
$Y\subseteq X_k$.  It follows that
\[
|X_k|\geq |Y|
\geq\frac{m_k}{\binom{n_k-1}{k-1}}
\geq\frac{\delta n}{s}.
\]
Thus in every case $|C(f(U))|\leq(1-\delta/s)n$, proving the first
implication.

For the converse, let $G=(V,E)$ be an $n$-vertex graph with $n\geq N$ that is
$\varepsilon$-far from $\Pi$.  The assertion is vacuous when $n<t$, so assume
$n\geq t$ and choose $T\in\binom Vt$ uniformly.  If $G[T]\in\Pi$, then for
some $j\leq q$ its fingerprint $F=f(T)$ is an ordered $j$-tuple contained in
$T$, and
\(T\subseteq C(F)\), where $|C(F)|\leq(1-\eta)n$.  For a fixed ordered tuple
$F$ of length $j$,
\(\Pr[F\subseteq T\subseteq C(F)] =\frac{\binom{|C(F)|-j}{t-j}}{\binom nt}.\)
There are at most $(n)_j$ possible ordered tuples of length $j$.  Summing over
all fingerprints and using
\(\frac{(n)_j}{\binom nt} =\frac{(t)_j}{\binom{n-j}{t-j}}\)
gives
\begin{align*}
\Pr[G[T]\in\Pi]
&\leq
\sum_{j=0}^{q}(t)_j
\frac{\binom{\lfloor(1-\eta)n\rfloor-j}{t-j}}
{\binom{n-j}{t-j}}\\
&\leq
\sum_{j=0}^{q}t^j(1-\eta)^{t-j}.
\end{align*}
Therefore at least a $\delta_t$-fraction of the $t$-subsets induce graphs
outside $\Pi$, as required.
\end{proof}

\end{document}
